% from ega0-22-fr.tex
\oldpage[0\textsubscript{IV-1}]{189}

En particulier, l'homomorphisme identique de $\Upsilon_{K/k}$ associe
\emph{canoniquement} à $K$ une $k$-algèbre noethérienne complète,
formellement lisse, pour laquelle $\Upsilon_{K/k}$ est isomorphe à
$\mathfrak m/\mathfrak m^2$.

\begin{remarks}[22.2.7]
\label{0.22.2.7-fr}
\begin{enumerate}
  \item[(i)] Soit $\Lambda$ un anneau local noethérien, de corps résiduel
  $k$~; il résulte de (19.7.1) et (19.7.2) que la détermination des
  $\Lambda$-algèbres $A$ qui sont des anneaux locaux noethériens complets et
  qui sont formellement lisses est \emph{équivalente} au même problème où
  $\Lambda$ est remplacé par $k$, et est donc en principe entièrement
  résolue par (22.2.5), qui ramène la question à la structure des \emph{corps
  résiduels} des $k$-algèbres cherchées, en tant qu'extensions de $k$.

  \item[(ii)] Le fait que $\Upsilon_{K/k}$ soit de rang fini sur $K$
  \emph{n'entraîne pas} que $K$ soit «~de multiplicité radicielle finie~» sur
  $k$ (cf. (19.6.6)). Par exemple, soient $k_0$ un corps parfait de
  caractéristique $p>0$, $k=k_0(T)$ le corps des fractions rationnelles à une
  indéterminée sur $k_0$ et $K=k^{p^{-\infty}}$ la plus petite extension
  parfaite de $k$. On a alors $\Omega^1_K=\Omega^1_{K/k}=0$ (21.4.4), donc
  par définition (20.6.1.1)
  $\Upsilon_{K/k}=\Omega^1_k\otimes_kK$, et comme $k^p=k_0(T^p)$, $T$ est
  une $p$-base de $k$ sur $k^p$, donc (21.4.5), $\Omega^1_k$ est de rang $1$
  sur $k$, et par suite $\Upsilon_{K/k}$ de rang $1$ sur $K$. Mais pour tout
  $n$ le corps résiduel de l'anneau local $K\otimes_kk^{p^{-n}}$ est
  isomorphe à $K$ donc n'est pas séparable sur $k^{p^{-n}}$.
\end{enumerate}
\end{remarks}

Le théorème (22.2.2) se généralise comme suit~:

\begin{proposition}[22.2.8]
\label{0.22.2.8-fr}
Sous les conditions de (22.2.1), supposons que la caractéristique $p$ de $k$
soit $>0$~; soit $(k_\alpha)$ une famille filtrante décroissante de
sous-corps de $k$, telle que
$\bigcap_\alpha k_\alpha(k^p)=k^p$. Alors la condition (ii) de (22.2.2) peut
se remplacer par l'une quelconque des suivantes~:
\begin{enumerate}
  \item[a)] Il existe $\alpha_0$ tel que l'homomorphisme canonique
  \[
    \Omega^1_{k/k_\alpha}\otimes_kK
    \longrightarrow\Omega^1_{A/k_\alpha}\otimes_AK
  \]
  soit injectif pour tout $\alpha\geqslant\alpha_0$.
  \item[b)] Il existe $\alpha_0$ tel que $A$ soit une $k$-algèbre
  formellement lisse (pour la topologie $\mathfrak m$-préadique) relativement
  à $k_\alpha$, pour tout $\alpha\geqslant\alpha_0$.
  \item[c)] Pour tout $\alpha$ et tout entier $n\geqslant0$,
  l'homomorphisme canonique
  \[
    \Omega^1_{k/k_\alpha}\otimes_k(A/\mathfrak m^{n+1})
    \longrightarrow
    \Omega^1_{A/k_\alpha}\otimes_A(A/\mathfrak m^{n+1})
  \]
  est inversible à gauche.
\end{enumerate}
\end{proposition}

\begin{proof}
En effet, si $A$ est une $k$-algèbre formellement lisse, elle l'est \emph{a
fortiori} relativement à tout sous-corps de $k$, ce qui entraîne $c)$ en
vertu de (20.7.2) et (19.1.7)~; il est clair que $c)$ implique $b)$, en vertu
de (20.7.2), et que $b)$ implique $a)$. Enfin, compte tenu de (22.2.3), il
reste à prouver que $a)$ entraîne que l'homomorphisme canonique
\[
  u:\Omega^1_k\otimes_kK\longrightarrow\Omega^1_A\otimes_AK
\]
est injectif. Or, on a pour tout $\alpha$ un diagramme commutatif
\[
\xymatrix@C=5em@R=3em{
  \Omega^1_k\otimes_kK \ar[r]^{u} \ar[d]_{v_\alpha} &
  \Omega^1_A\otimes_AK \ar[d] \\
  \Omega^1_{k/k_\alpha}\otimes_kK \ar[r]_{u_\alpha} &
  \Omega^1_{A/k_\alpha}\otimes_AK
}
\]

